% Options for packages loaded elsewhere
% Options for packages loaded elsewhere
\PassOptionsToPackage{unicode}{hyperref}
\PassOptionsToPackage{hyphens}{url}
\PassOptionsToPackage{dvipsnames,svgnames,x11names}{xcolor}
%
\documentclass[
  letterpaper,
  DIV=11,
  numbers=noendperiod]{scrartcl}
\usepackage{xcolor}
\usepackage{amsmath,amssymb}
\setcounter{secnumdepth}{-\maxdimen} % remove section numbering
\usepackage{iftex}
\ifPDFTeX
  \usepackage[T1]{fontenc}
  \usepackage[utf8]{inputenc}
  \usepackage{textcomp} % provide euro and other symbols
\else % if luatex or xetex
  \usepackage{unicode-math} % this also loads fontspec
  \defaultfontfeatures{Scale=MatchLowercase}
  \defaultfontfeatures[\rmfamily]{Ligatures=TeX,Scale=1}
\fi
\usepackage{lmodern}
\ifPDFTeX\else
  % xetex/luatex font selection
\fi
% Use upquote if available, for straight quotes in verbatim environments
\IfFileExists{upquote.sty}{\usepackage{upquote}}{}
\IfFileExists{microtype.sty}{% use microtype if available
  \usepackage[]{microtype}
  \UseMicrotypeSet[protrusion]{basicmath} % disable protrusion for tt fonts
}{}
\makeatletter
\@ifundefined{KOMAClassName}{% if non-KOMA class
  \IfFileExists{parskip.sty}{%
    \usepackage{parskip}
  }{% else
    \setlength{\parindent}{0pt}
    \setlength{\parskip}{6pt plus 2pt minus 1pt}}
}{% if KOMA class
  \KOMAoptions{parskip=half}}
\makeatother
% Make \paragraph and \subparagraph free-standing
\makeatletter
\ifx\paragraph\undefined\else
  \let\oldparagraph\paragraph
  \renewcommand{\paragraph}{
    \@ifstar
      \xxxParagraphStar
      \xxxParagraphNoStar
  }
  \newcommand{\xxxParagraphStar}[1]{\oldparagraph*{#1}\mbox{}}
  \newcommand{\xxxParagraphNoStar}[1]{\oldparagraph{#1}\mbox{}}
\fi
\ifx\subparagraph\undefined\else
  \let\oldsubparagraph\subparagraph
  \renewcommand{\subparagraph}{
    \@ifstar
      \xxxSubParagraphStar
      \xxxSubParagraphNoStar
  }
  \newcommand{\xxxSubParagraphStar}[1]{\oldsubparagraph*{#1}\mbox{}}
  \newcommand{\xxxSubParagraphNoStar}[1]{\oldsubparagraph{#1}\mbox{}}
\fi
\makeatother


\usepackage{longtable,booktabs,array}
\usepackage{calc} % for calculating minipage widths
% Correct order of tables after \paragraph or \subparagraph
\usepackage{etoolbox}
\makeatletter
\patchcmd\longtable{\par}{\if@noskipsec\mbox{}\fi\par}{}{}
\makeatother
% Allow footnotes in longtable head/foot
\IfFileExists{footnotehyper.sty}{\usepackage{footnotehyper}}{\usepackage{footnote}}
\makesavenoteenv{longtable}
\usepackage{graphicx}
\makeatletter
\newsavebox\pandoc@box
\newcommand*\pandocbounded[1]{% scales image to fit in text height/width
  \sbox\pandoc@box{#1}%
  \Gscale@div\@tempa{\textheight}{\dimexpr\ht\pandoc@box+\dp\pandoc@box\relax}%
  \Gscale@div\@tempb{\linewidth}{\wd\pandoc@box}%
  \ifdim\@tempb\p@<\@tempa\p@\let\@tempa\@tempb\fi% select the smaller of both
  \ifdim\@tempa\p@<\p@\scalebox{\@tempa}{\usebox\pandoc@box}%
  \else\usebox{\pandoc@box}%
  \fi%
}
% Set default figure placement to htbp
\def\fps@figure{htbp}
\makeatother





\setlength{\emergencystretch}{3em} % prevent overfull lines

\providecommand{\tightlist}{%
  \setlength{\itemsep}{0pt}\setlength{\parskip}{0pt}}



 


\usepackage{booktabs}
\usepackage{longtable}
\usepackage{array}
\usepackage{multirow}
\usepackage{wrapfig}
\usepackage{float}
\usepackage{colortbl}
\usepackage{pdflscape}
\usepackage{tabu}
\usepackage{threeparttable}
\usepackage{threeparttablex}
\usepackage[normalem]{ulem}
\usepackage{makecell}
\usepackage{xcolor}
\usepackage{caption}
\usepackage{anyfontsize}
\KOMAoption{captions}{tableheading}
\makeatletter
\@ifpackageloaded{caption}{}{\usepackage{caption}}
\AtBeginDocument{%
\ifdefined\contentsname
  \renewcommand*\contentsname{Table of contents}
\else
  \newcommand\contentsname{Table of contents}
\fi
\ifdefined\listfigurename
  \renewcommand*\listfigurename{List of Figures}
\else
  \newcommand\listfigurename{List of Figures}
\fi
\ifdefined\listtablename
  \renewcommand*\listtablename{List of Tables}
\else
  \newcommand\listtablename{List of Tables}
\fi
\ifdefined\figurename
  \renewcommand*\figurename{Figure}
\else
  \newcommand\figurename{Figure}
\fi
\ifdefined\tablename
  \renewcommand*\tablename{Table}
\else
  \newcommand\tablename{Table}
\fi
}
\@ifpackageloaded{float}{}{\usepackage{float}}
\floatstyle{ruled}
\@ifundefined{c@chapter}{\newfloat{codelisting}{h}{lop}}{\newfloat{codelisting}{h}{lop}[chapter]}
\floatname{codelisting}{Listing}
\newcommand*\listoflistings{\listof{codelisting}{List of Listings}}
\makeatother
\makeatletter
\makeatother
\makeatletter
\@ifpackageloaded{caption}{}{\usepackage{caption}}
\@ifpackageloaded{subcaption}{}{\usepackage{subcaption}}
\makeatother
\usepackage{bookmark}
\IfFileExists{xurl.sty}{\usepackage{xurl}}{} % add URL line breaks if available
\urlstyle{same}
\hypersetup{
  pdftitle={PIRL data exercise report},
  pdfauthor={Martin Chen (bc3517@nyu.edu)},
  colorlinks=true,
  linkcolor={blue},
  filecolor={Maroon},
  citecolor={Blue},
  urlcolor={Blue},
  pdfcreator={LaTeX via pandoc}}


\title{PIRL data exercise report}
\author{Martin Chen (bc3517@nyu.edu)}
\date{2026-02-23}
\begin{document}
\maketitle

\renewcommand*\contentsname{Table of contents}
{
\hypersetup{linkcolor=}
\setcounter{tocdepth}{3}
\tableofcontents
}

\newpage{}

\section{Introduction}\label{introduction}

\subsection{Structure of the analysis
folder}\label{structure-of-the-analysis-folder}

\begin{itemize}
\item
  \texttt{\_targets}

  \begin{itemize}
  \tightlist
  \item
    A folder containing all the outputs from \texttt{\_targets.R}. It
    can only be loaded to the R environment using the \texttt{targets}
    package in R
  \item
    If you want to remove this folder and recreate another one, please
    set the working directory to the location of \texttt{\_targets.R},
    execute \texttt{targets::tar\_destroy()}, and then run
    \texttt{targets::tar\_make()}
  \end{itemize}
\item
  \texttt{\_targets.R}

  \begin{itemize}
  \item
    An \texttt{.R} script to reproduce the analysis. To reproduce the
    analysis, you can set the working directory to the location of this
    script and execute \texttt{targets::tar\_make()} on the console.
  \item
    If you don't want to reproduce the results, you can just set the
    working directory to the location of this script and execute the R
    command below on the console to check the established results.
  \end{itemize}
\end{itemize}

\begin{verbatim}
if (!require('targets', character.only = TRUE)) {
  install.packages('targets', dependencies = TRUE)
  library('targets')
}

library(targets)

# cleaed dataset for analysis
tar_load(us_datasets)
# logistic regression outcome of changed_behaviors ~ outside_spaces
tar_load(result_logistic)
# logistic regression outcome of changed_behaviors ~ outside_spaces + income
tar_load(result_logistic_income)
# chi-suqared test results of changed_behaviors * outside_spaces
tar_load(chisq_bc_os)
# poissson regression (w/ robust standard error) outcome of changed_behaviors ~ outside_spaces
tar_load(result_poisson_RSE)
# probit regression outcome of changed_behaviors ~ outside_spaces
tar_load(result_probit)
\end{verbatim}

\begin{itemize}
\item
  \texttt{R}

  \begin{itemize}
  \tightlist
  \item
    A folder containing all the \texttt{R} functions used in the
    \texttt{\_targets.R} script.
  \end{itemize}
\item
  \texttt{raw\_data}

  \begin{itemize}
  \tightlist
  \item
    Raw datasets and codebooks for this data exercise.
  \end{itemize}
\end{itemize}

\section{Question 1 \& 2: Data input and
cleaning}\label{question-1-2-data-input-and-cleaning}

\subsection{Data input}\label{data-input}

\begin{itemize}
\tightlist
\item
  file location: ``raw\_data/Data\_Exercise\_CSV.csv''
\end{itemize}

\subsection{Data cleaning}\label{data-cleaning}

\begin{itemize}
\tightlist
\item
  Martin chose the country ``US.''

  \begin{itemize}
  \tightlist
  \item
    After the selection, the dataset had 1055 observations.
  \item
    Please refer to for Table~\ref{tbl-of-us-summary} summary statistics
  \end{itemize}
\end{itemize}

\begingroup
\fontsize{12.0pt}{14.0pt}\selectfont
\setlength{\LTpost}{0mm}

\begin{longtable}{lc}

\caption{\label{tbl-of-us-summary}Summary statistics for observation in
the U.S.}

\tabularnewline

\toprule
 & \multicolumn{1}{c}{{U.S.}} \\ 
\cmidrule(lr){2-2}
\textbf{Characteristic} & \textbf{N = 1,055}\textsuperscript{\textit{1}} \\ 
\midrule\endhead\addlinespace[2.5pt]
Changes in behaviors &  \\ 
    No & 128 (12\%) \\ 
    Yes & 927 (88\%) \\ 
Income quantiles of participants &  \\ 
    First quintile & 181 (17\%) \\ 
    Second quintile & 196 (19\%) \\ 
    Third quintile & 218 (21\%) \\ 
    Fourth quintile & 245 (23\%) \\ 
    Fifth quintile & 199 (19\%) \\ 
    Prefer not to answer & 16 (1.5\%) \\ 
Access to outside space at home &  \\ 
    No & 162 (15\%) \\ 
    Yes & 893 (85\%) \\ 
\bottomrule

\end{longtable}

\begin{minipage}{\linewidth}
\vspace{.05em}
\textsuperscript{\textit{1}}n (\%)\\
\end{minipage}
\endgroup

\section{Question 3: Regression on behavior change and access to outside
space}\label{question-3-regression-on-behavior-change-and-access-to-outside-space}

\subsection{Model used}\label{model-used}

\begin{itemize}
\item
  Martin chose a logistic regression model to estimate the relationship
  between behavior change and access to outside space.

  \[
  log(\frac{P_{behavior\ change = 1}}{1- P_{behavior\ change = 1}}) \sim access\ to\ outside\ space
  \]
\item
  Regression result

  \begin{itemize}
  \tightlist
  \item
    Please refer to Table~\ref{tbl-of-regression-single} for detailed
    regression results
  \end{itemize}
\end{itemize}

\begin{table}

\caption{\label{tbl-of-regression-single}Regression result of question
3}

\centering{

\fontsize{12.0pt}{14.0pt}\selectfont
\begin{tabular*}{\linewidth}{@{\extracolsep{\fill}}lcccccc}
\toprule
 & \multicolumn{3}{c}{{\textbf{Univariate Model}}} & \multicolumn{3}{c}{{\textbf{Income added}}} \\ 
\cmidrule(lr){2-4} \cmidrule(lr){5-7}
\textbf{Characteristic} & \textbf{OR} & \textbf{95\% CI} & \textbf{p-value} & \textbf{OR} & \textbf{95\% CI} & \textbf{p-value} \\ 
\midrule\addlinespace[2.5pt]
Access to outside space at home &  &  &  &  &  &  \\ 
    No & 1.00 & — &  & 1.00 & — &  \\ 
    Yes & 1.94 & 1.23, 3.00 & 0.003 & 1.59 & 0.99, 2.52 & 0.049 \\ 
Income quantile &  &  &  &  &  &  \\ 
    First quintile &  &  &  & 1.00 & — &  \\ 
    Second quintile &  &  &  & 0.83 & 0.49, 1.40 & 0.5 \\ 
    Third quintile &  &  &  & 1.57 & 0.88, 2.83 & 0.13 \\ 
    Fourth quintile &  &  &  & 2.26 & 1.23, 4.26 & 0.010 \\ 
    Fifth quintile &  &  &  & 3.62 & 1.76, 8.06 & <0.001 \\ 
    Prefer not to answer &  &  &  & 0.43 & 0.14, 1.44 & 0.14 \\ 
\bottomrule
\end{tabular*}
\begin{minipage}{\linewidth}
\vspace{.05em}
Abbreviations: CI = Confidence Interval, OR = Odds Ratio\\
\end{minipage}

}

\end{table}%

\subsection{Alternative models}\label{alternative-models}

\begin{itemize}
\tightlist
\item
  Relative to probit, logit has a practical advantage:

  \begin{itemize}
  \tightlist
  \item
    Logit coefficients have a log-odds interpretation, which is
    sometimes more intuitive.
  \end{itemize}
\item
  Relative to Poisson regression with robust standard errors

  \begin{itemize}
  \tightlist
  \item
    The outcome is not a count, so the Poisson likelihood is
    misspecified.
  \item
    Interpretation is less intuitive than the logit model.
  \end{itemize}
\item
  Compared with the raw answer to outside space access

  \begin{itemize}
  \tightlist
  \item
    In the original survey questionnaire, the answer to the question
    about outside space included 7 choices.
  \item
    Conversion to a binary variable makes it conceptually simple and
    policy-relevant: does having any access to outside space matter,
    relative to having none?
  \item
    The conversion also simplifies model interpretation.
  \end{itemize}
\end{itemize}

\section{Question 4}\label{question-4}

\subsection{Model interpretation}\label{model-interpretation}

\begin{itemize}
\tightlist
\item
  Compared with participants without access to the outside space, those
  with access had a higher odds of changing their behavior by a factor
  of 1.94 on average in our United States samples.
\end{itemize}

\subsection{Discussion on causal
interpretation}\label{discussion-on-causal-interpretation}

\begin{itemize}
\item
  Given the model design and the dataset's background, Martin did not
  find sufficient evidence to support a causal interpretation.
\item
  \texttt{outside\_space} is not randomly assigned. Instead, it is an
  outcome of participants' housing choices and constraints, shaped by
  various factors. Therefore, the estimated coefficient reflected the
  effect of outdoor space and the influence of omitted variables related
  to it.
\end{itemize}

\subsection{Omitted variables}\label{omitted-variables}

\begin{itemize}
\item
  A range of omitted variables could be correlated with outside access.

  \begin{itemize}
  \item
    Housing type (houses versus apartments), urban density, neighborhood
    characteristics, household size, presence of children, education,
    ability to work from home, local COVID-19 prevalence, and local
    policy stringency are all plausibly correlated with both access to
    outside space and behavior change
  \item
    Omitting these variables could lead to an upward bias in the
    coefficient for outside space, as shown in the model with income
    quantiles.
  \end{itemize}
\end{itemize}

\subsection{Endogenity}\label{endogenity}

\begin{itemize}
\tightlist
\item
  Outside space is endogenous because it is determined by residential
  location, housing quality, and household resources, all of which may
  affect behavior change.
\end{itemize}

\subsection{Ideal experiment}\label{ideal-experiment}

\begin{itemize}
\item
  An ideal experiment to identify the causal effect of outside space
  would randomly assign individuals to otherwise identical housing units
  that differ only in access to outside space, and then measure
  subsequent behavior change.
\item
  This could be carried out through randomized housing lotteries or
  quasi-random assignment of public or subsidized housing units.
\item
  In the current analysis, access to outside space was non-random and
  observational. Martin thought some confounding factors were
  unobserved.
\end{itemize}

\subsection{Threats to identification}\label{threats-to-identification}

\begin{itemize}
\item
  Several threats to identification are non-random residential sorting,
  unobserved risk preferences, differential exposure to information, and
  local policy variation correlated with housing characteristics.

  \begin{itemize}
  \item
    One solution is to add more control variables, including detailed
    housing characteristics, neighborhood fixed effects, local COVID-19
    case rates, and policy measures.
  \item
    Another is to use individual panel data to allow for individual
    fixed effects.
  \end{itemize}
\end{itemize}

\subsection{Direction of bias}\label{direction-of-bias}

\begin{itemize}
\item
  Martin believes the bias could be upwards.

  \begin{itemize}
  \tightlist
  \item
    Failure to control other confounders may exaggerate the effects of
    outside
  \end{itemize}
\end{itemize}

\subsection{Possible use of instrument variables
(IV)}\label{possible-use-of-instrument-variables-iv}

\begin{itemize}
\item
  A possible IV is waitlist-based allocation of units with and without
  outdoor space, such as affordable housing policies in cities like NYC
  and Houston.
\item
  The ideal IV should only affect behavior change through access to
  outside space and not through other ways.
\end{itemize}

\section{Question 5}\label{question-5}

\subsection{Model interpretation}\label{model-interpretation-1}

\begin{itemize}
\tightlist
\item
  Compared with participants without access to the outside, those with
  access had a 1.59-fold increase in the odds of changing their
  behavior, on average, in our United States samples, controlling for
  all other covariates.
\end{itemize}

\subsection{Differences from step 3}\label{differences-from-step-3}

\begin{itemize}
\item
  The odds ratio for \texttt{outside\ space} was lower, indicating that
  its effect was weaker when income quantile was controlled for in the
  model.
\item
  The p-value for \texttt{outside\ space} was higher, indicating that
  Martin had weaker evidence against the null hypothesis that the
  outside space has no effect on the behavior change.
\item
  Including income directly addresses one important source of
  endogeneity: selection into housing with outside space based on
  socioeconomic status. However, the estimate should still be
  interpreted as correlational rather than causal.
\end{itemize}

\section{Question 6}\label{question-6}

\subsection{Policy Implications}\label{policy-implications}

\begin{itemize}
\item
  In the short run, policies should focus on expanding effective access
  to safe outdoor space, particularly for households living in dense or
  space-constrained housing. Potential interventions include keeping
  parks, playgrounds, and open streets accessible during public health
  emergencies; creating temporary outdoor areas near high-density
  housing (e.g., pop-up pedestrian zones or shared courtyards); and
  providing clear guidance to building managers to keep shared outdoor
  spaces open and safely managed rather than fully closed. These
  policies may be effective because access to outdoor space can reduce
  the psychological and practical costs of staying home, making
  compliance with broader behavioral guidelines more feasible.
\item
  In the long run, the findings point to housing and urban planning
  policies that incorporate outdoor access as a public health input.
  Examples include encouraging or subsidizing the construction of
  balconies, terraces, or ventilated shared courtyards in multi-family
  housing; revising zoning or building codes to prioritize green space
  in high-density areas; and targeting housing improvements toward
  lower-income neighborhoods that are less likely to have private
  outdoor access. Such policies may have lasting benefits beyond
  pandemics by improving mental health, well-being, and resilience to
  future public health shocks.
\end{itemize}

\subsection{Additional Quantitative Data
Collection}\label{additional-quantitative-data-collection}

\begin{itemize}
\item
  First, housing characteristics would be critical. Data on housing type
  (houses or apartments), number of rooms, crowding (rooms per person),
  ownership status, and whether outdoor space is private or shared would
  allow for a more precise characterization of living conditions and
  help distinguish the mechanisms through which outside space operates.
\item
  Second, individual and household characteristics such as age,
  education, employment status, ability to work from home, household
  size, presence of children or elderly members, and baseline health
  risk would help account for differences in both exposure risk and
  capacity to change behavior. These variables would reduce
  omitted-variable bias and enable heterogeneity analyses across
  demographic groups.
\item
  Third, local context variables, such as neighborhood density, urban
  versus suburban location, local COVID-19 case rates, and policy
  stringency measures, would help separate the effect of outside space
  from broader environmental and policy conditions that influence
  behavior.
\item
  Finally, behavioral detail and dynamics would improve interpretation.
  Rather than a single indicator for behavior change, more granular
  measures (e.g., masking, social distancing, avoidance of gatherings,
  time spent at home) and panel data tracking individuals over time
  would allow researchers to examine which behaviors are most affected
  by outside space and whether effects persist or evolve as conditions
  change.
\end{itemize}

\section{Question 7}\label{question-7}

\subsection{Methods}\label{methods}

\begin{itemize}
\tightlist
\item
  I would collect qualitative data through in-depth, semi-structured
  interviews, potentially supplemented by ethnographic observation.
\end{itemize}

\subsection{Questions asked}\label{questions-asked}

\begin{itemize}
\item
  I would conduct open-ended interviews with individuals or households,
  using a semi-structured interview guide.
\item
  First, I would ask respondents to describe how their daily behaviors
  changed during the pandemic, including specific behaviors such as
  staying home, avoiding gatherings, mask use, and use of public spaces.

  \begin{itemize}
  \tightlist
  \item
    This helps clarify what ``behavior change'' means in practice and
    whether different respondents interpret it similarly.
  \end{itemize}
\item
  Second, I would ask detailed questions about housing and outside space

  \begin{itemize}
  \item
    What types of outside space they had access to (private, shared,
    public nearby)
  \item
    How frequently and for what purposes they used these spaces
  \item
    Whether outside space affected their willingness or ability to
    comply with public health guidance
  \item
    These questions are designed to uncover the mechanisms through which
    outside space may facilitate behavior change
  \end{itemize}
\item
  Third, I would investigate constraints and trade-offs, asking
  respondents what made behavior change difficult, what they felt they
  were giving up, and how housing conditions shaped those trade-offs.

  \begin{itemize}
  \tightlist
  \item
    This can reveal barriers that are difficult to quantify, such as
    mental health strain, caregiving responsibilities, or household
    conflicts.
  \end{itemize}
\item
  Finally, I would ask respondents how local rules, building management
  decisions, and neighborhood norms influenced their behavior.

  \begin{itemize}
  \tightlist
  \item
    These questions help contextualize individual decisions within
    broader institutional and social environments.
  \end{itemize}
\end{itemize}

\subsection{Sampling strategy}\label{sampling-strategy}

\begin{itemize}
\tightlist
\item
  I would use purposeful sampling, focusing on respondents with key
  characteristics (e.g., high versus low income levels, single apartment
  or family house, urban versus suburban areas)
\end{itemize}

\subsection{How qualitative data augments the quantitative
analysis}\label{how-qualitative-data-augments-the-quantitative-analysis}

\begin{itemize}
\item
  First, qualitative analysis can interpret the regression coefficients
  by explaining how behavior changed in real life and how outside space
  affected it.
\item
  Second, qualitative analysis can help identify omitted variables, such
  as social norms, building rules, or conditions of personal well-being.
\end{itemize}

\subsection{Ethnographic observation}\label{ethnographic-observation}

\begin{itemize}
\tightlist
\item
  Yes. I would incorporate an ethnographic observation, such as
  observing public spaces in apartment buildings and neighborhoods.
\end{itemize}

\subsection{Values of ethnographic
observation}\label{values-of-ethnographic-observation}

\begin{itemize}
\item
  Ethnographic observations could document the condition of outdoor
  spaces, how they were actually used, and the nature of social
  interactions.
\item
  These observations could help researchers identify factors respondents
  ignored, such as their actual behaviors rather than their stated ones.
\end{itemize}

\subsection{Effects on answers to step
6}\label{effects-on-answers-to-step-6}

\begin{itemize}
\item
  The observations could help understand the mechanisms and reasons
  behind the effects of the external environment on behavior change.
\item
  For example, simply expanding nominal access to public spaces may be
  insufficient and could even be counterproductive if those spaces are
  perceived as unsafe or become associated with disorder. Effective
  policy design therefore requires ensuring that outdoor spaces are
  safe, well-managed, and socially acceptable in practice.
\end{itemize}




\end{document}
